To fortify our intuitions, we examine a concrete mathematical example in detail.

We start out with a theory \expr{PosPlus} of positive natural numbers with addition
and intersect it with \expr{StrConc} of strings with concatenation (as in \autoref{fig:posstr}). Note that we do
not start with modular developments; rather the modular structure is (going to be) the
result of intersecting with various examples. Also, note that the views \lstinline|k| and \lstinline|l| are both partial.

\begin{figure}[ht]\centering
\begin{tabular}{|p{6cm}|p{6cm}|}\hline
  Positive  & Strings \\\hline
\begin{mmtcode}[aboveskip=-2pt,belowskip=-2pt]
theory PosPlus =
    ℕ : type ❙
    + : ℕ ⟶ ℕ ⟶ ℕ
    assoc : ...
    comm : ...
\end{mmtcode}
& 
\begin{mmtcode}[aboveskip=-2pt,belowskip=-2pt]
theory StrConc =
    A* : type ❙
    :: : A* ⟶ A* ⟶ A* ❙
    assoc : ...
    ε : A* ❙
    unit : ... ❙
\end{mmtcode}
\\[-2pt]\hline
\begin{mmtcode}[aboveskip=-2pt,belowskip=-2pt]
view k : ?PosPlus -> ?StrConc =
    ℕ = A* ❙
    + = :: ❙
    assoc = assoc ❙
\end{mmtcode}
&
\begin{mmtcode}[aboveskip=-2pt,belowskip=-2pt]
view l : ?StrConc -> ?PosPlus =
    A* = ℕ ❙
    :: = + ❙
    assoc = assoc ❙
\end{mmtcode}
\\[-2pt]\hline
\end{tabular}
\caption{Positive Integers and Strings}\label{fig:posstr}\vspace*{-1.5em}
\end{figure}

Now the intersection as proposed above yields
\autoref{fig:posstr-intersection}, which directly corresponds to the schema in
\autoref{fig:theory-intersection}. 
Note, that the new pair of (equivalent) theories is completely
determined by the partial morphism \cn{k}; the only thing we have to invent are the names
-- and we have not been very inventive here.

\begin{figure}[ht]\centering
\begin{tabular}{|p{6cm}|p{6cm}|}\hline
  Positive& Strings \\[-2pt]\hline
\begin{mmtcode}[aboveskip=-2pt,belowskip=-2pt]
theory A =
    ℕ : type ❙
    + : ℕ ⟶ ℕ ⟶ ℕ ❙
    assoc : ... 
❚
theory PosPlus =
    include ?A ❙
    comm : ...
❚
\end{mmtcode}
& 
\begin{mmtcode}[aboveskip=-2pt,belowskip=-2pt]
theory B =
    A* : type ❙
    :: : A* ⟶ A* ⟶ A* ❙
    assoc : ... ❙
❚
theory StrConc =
    include ?B ❙
    ε : A* ❙
    unit : ... ❙
❚
\end{mmtcode}
\\[-2pt]\hline
\end{tabular}
\caption{Positive Integers and Strings, Refactored}\label{fig:posstr-intersection}\vspace*{-1.5em}
\end{figure}

Intuitively, the intersection theory is the theory of semigroups which is traditionally
written down in \mmt as in \autoref{fig:sg}. And indeed, \expr{sg} is a
renaming of both \expr{A} and \expr{B}. 

\begin{figure}[ht]%{r}{5.3cm}\vspace*{-2.8em}\hspace*{-1em}
\centering
\begin{lstlisting}
theory sg =
    G : type ❙
    ∘ : G ⟶ G ⟶ G ❙
\end{lstlisting}%\vspace*{-1em}
\caption{Semigroups}\label{fig:sg}%\vspace*{-2.5em}
\end{figure}

For this situation, we should have a variant theory
intersection operator that takes a partial morphism and returns \emph{one} intersection theory.
   

In this situation, we want a theory intersection operator that follows the schema in \autoref{fig:unary-thi}. 

\begin{figure}[ht]%l{4.2cm}\vspace*{-2em}\hspace*{-1em}
\centering
\tikzinput{img/int-unary-thi}%\vspace*{-.5em}
\caption{Unary TI}\label{fig:unary-thi}%\vspace*{-.5em}
\end{figure}

Let us call this operation \emph{unary TI} and the previous one \emph{binary TI}. To compute it, we need to specify a name $N$ for the new theory and two renamings
$\rho_1:\viewto N{\Domain\sigma}$ and
$\rho_2:\viewto N{\textsf{Img}(\sigma)}$. Note that in
this TI operator, the intersection is connected to the ``rest theories'' via \mmt
structures -- rather than mere inclusions -- which carry the assignments induced by the
partial morphisms suitably composed with the renamings (see next section).

In our example we can obtain the theory \cn{sg} from \autoref{fig:sg} via the renamings
$\rho_1:=\{\cn{G}\mmtass\mathbb{N},\circ\mmtass+,\cn{assoc}\mmtass\cn{assoc}\}$ and
$\rho_2:=\{\cn{G}\mmtass A^*,\circ\mmtass\conc,\cn{assoc}\mmtass\cn{assoc}\}$.

Unary TI is often the more useful operation on theory graphs, but needs direct user
supervision, whereas binary TI is fully automatic if we accept generated names for the
intersection theories.
